%% kapa: exam=1 round=28 year=2017 period=2 form=A subject=accounting subject_ko=회계학 seq=2 number=42
%% source: Q-Net 감정평가사 기출문제 (cst00309) artlSeq=5199603 file=extracted/1008/5199603/2017년 제28회 감정평가사 1차시험 문제지/2교시 A형.pdf
%% conversion: kordoc 4.24.51 -> tools/gen1.py (auto); review=proofread-sonnet; flags=none
\begin{question}{42}
수익인식에 관한 설명으로 옳지 않은 것은?
\begin{choices}
  \item 재화의 결함에 대하여 정상적인 품질보증범위를 초과하여 판매자가 책임을 지더라도 재화가 구매자에게 인도되었다면 수익으로 인식한다.
  \item 설치조건부 판매에서 계약의 중요한 부분을 차지하는 설치가 아직 완료되지 않은 경우에는 당해 거래를 판매로 보지 아니하며, 수익을 인식하지 아니한다.
  \item 위탁판매의 경우 위탁자는 수탁자가 제3자에게 재화를 판매한 시점에 수익을 인식한다.
  \item 주문형 소프트웨어의 개발수수료는 진행기준에 따라 수익을 인식한다.
  \item 판매대금의 회수가 구매자의 재판매에 의해 결정되는 경우에는 당해 거래를 판매로 보지 아니하며, 수익을 인식하지 아니한다.
\end{choices}
\end{question}
\answer{1}
